\section{Conclusion}
\label{sec:conclusion}

This paper presented a controlled diagnostic study of HyDE-style document-like generative query expansion in a fixed-depth, training-free protocol for multi-hop RAG. Across HotpotQA and the secondary 2WikiMultihopQA check, the main observed pattern is better complete supporting-title retrieval, with matching paragraph-text audit results and reader EM/F1 gains reported for the fixed reader used in the end-to-end pipeline. The matched-call query-reformulation controls show that the retrieval gain is not explained by an extra query-side LLM call alone. The strict equal-budget retrieval-only diagnostic further shows that independently generated document-like passages outperform direct rewrites on complete supporting-title retrieval at closely matched generated lengths and call budgets. The secondary guarded canonicalization diagnostic remains an answer-format risk analysis rather than part of the main empirical claim.

Overall, the evidence supports a retrieval-stage conclusion: under the evaluated processed local-corpus protocol, document-like hypothetical expansion improves supporting-title and paragraph-text acquisition over direct reformulation and question decomposition at a similar generation budget, while its comparison with keyword/entity expansion is dataset-dependent. The targeted Qwen-Turbo readout is consistent with this boundary rather than a broad reader-independent superiority claim. Sentence-level supporting-fact coverage remains outside the available processed artifacts, and code/artifact availability is described in Supplement Section~\ref{supp:sec:supp_code_artifact_availability}.
